\documentclass[10pt,a4paper]{amsart}
\usepackage[margin=19mm]{geometry}
\usepackage{amsmath,amssymb,mathtools}
\usepackage[scr=rsfs,cal=euler]{mathalfa}
\usepackage{xfrac}
\usepackage[unicode,hidelinks,bookmarks=false]{hyperref}
\newcommand{\ie}{i.e.\ }
\newcommand{\cf}{cf.\ }
\newcommand{\resp}{respectively\ }
\newcommand{\Subsection}{\subsection}
\providecommand{\nobreakdash}{\nobreak\hbox{-}}
\setlength{\emergencystretch}{2em}
\multlinegap0pt
\usepackage{amssymb}
\usepackage{tikz}
\newtheorem{theorem}{Theorem}
\newtheorem{proposition}{Proposition}
\title{Binary LCD Codes and Their Graph Representations}
\author{Keita Ishizuka}
\date{Source: arXiv:2407.07689v5 (CC BY 4.0)}
\begin{document}
\maketitle

\noindent\emph{Source and licence.} Binary LCD Codes and Their Graph Representations. Version arXiv:2407.07689v5 (CC BY 4.0), distributed by arXiv under the Creative Commons Attribution 4.0 International licence, http://creativecommons.org/licenses/by/4.0/, as recorded in the arXiv licence metadata for this exact version and shown on the abstract page https://arxiv.org/abs/2407.07689v5. Full licence notice: \texttt{licenses/arxiv-2407.07689-CC-BY-4.0.txt}.\par\smallskip

\noindent\textit{Editorial scope.} Counted unit: the article's proposition proposition (source0.tex lines 375--377). The article's complete original proof of it is delivered with it (source0.tex lines 378--412, one \texttt{proof} environment of 35 lines). No mathematics is written, repaired or re-derived: the statement and the proof are the article's own lines, in the article's own notation, and no retained line is substituted or re-flowed.  Retained uncounted as the source-local context the proof needs: lines 279--281.  Nothing was omitted from any retained span, so the package carries no omission record.  No retained line contains a citation, so the wrapper carries no bibliography and no bibliography entry.  No retained line contains a figure environment, an image-inclusion command or drawing code, so no image is delivered and the package declares no asset dependency.  Editorial formatting only, in addition: an AMS \texttt{amsart} preamble that restates the article's own declarations and macros used by the retained lines, namely the environments \texttt{theorem}, \texttt{proposition} on separate counters, together with the standard packages those lines need.  The retained environments print in this document as Theorem 1, Proposition 1 in order of appearance, whereas the article prints the same items with the numbering of its own full document.  The complete native source of the article as distributed in the arXiv e-print for this exact version is bundled unchanged as \texttt{sources/162119/source0.tex}.
\begin{theorem}[DRG Characterization]\label{thm:drg-main}
A distance-regular graph $\Gamma$ with intersection array $\{b_0, b_1, \ldots, b_{d-1}; c_1, c_2, \ldots, c_d\}$ yields a binary even LCD code if and only if (i) $b_0 \equiv 0 \pmod 2$; (ii) $a_1 \equiv 1 \pmod 2$ where $a_1 = k - b_1 - c_1$; and (iii) $c_2 \equiv 0 \pmod 2$.
\end{theorem}

\begin{proposition}[Grassmann Graph Characterization]
The Grassmann graph $J_q(n,k)$ with $q$ odd yields a binary even LCD code if and only if $n$ is odd. If $q$ is even, then $J_q(n,k)$ never yields a binary even LCD code.
\end{proposition}

\begin{proof}
The Grassmann graph $J_q(n,k)$ has intersection array with
$$
b_i = q^{2i+1}\binom{k-i}{1}_q \binom{n-k-i}{1}_q,\ c_i = \binom{i}{1}_q^2.
$$
A straightforward calculation shows that
\begin{align*}
b_i
&= q^{2i+1} \binom{k-i}{1}_q \binom{n-k-i}{1}_q\\
&= \frac{q^{2i+1}(q^{k-i}-1)(q^{n-k-i}-1)}{(q-1)(q-1)}\\
&= q^{2i+1} (q^{k-i-1}+q^{k-i-2}+\dots+1)(q^{n-k-i-1}+q^{n-k-i-2}+\dots+1)\\
&\equiv (k-i)(n-k-i) \pmod 2.
\end{align*}
In the fourth equality, we used that $q \equiv 1 \pmod 2$.
In a similar manner, we see that
$$
c_i
= \binom{i}{1}_q^2
= (\frac{q^i-1}{q-1})^2
= (q^{i-1}+ \dots + q+1)^2
\equiv i \pmod 2.
$$
Using these equations, we obtain
\begin{align*}   
a_1
&= k - b_1 - c_1\\
&\equiv k(n-k) - (k-1)(n-k-1) - 1 \pmod 2\\
&\equiv kn - k^2 - (kn - k^2 - k - n + k + 1) - 1 \pmod 2\\
&\equiv n \pmod 2.
\end{align*}
By Theorem~\ref{thm:drg-main}, $J_q(n,k)$ yields an LCD code if and only if $b_0 \equiv 0 \pmod 2$, $a_1 \equiv 1 \pmod 2$, and $c_2 \equiv 0 \pmod 2$. With the above reductions, these become $k(n-k) \equiv 0$, $n \equiv 1$, and $c_2 \equiv 0 \pmod 2$, respectively. Since $c_2 \equiv 0$ is automatically satisfied, the conditions reduce to $n$ being odd.

If $q$ is even, the above reductions do not apply.
Instead, since $q \equiv 0 \pmod 2$, we have $q^j \equiv 0 \pmod 2$ for $j \geq 1$, so $\binom{r}{1}_q = q^{r-1}+\cdots+q+1 \equiv 1 \pmod 2$ for all $r \geq 1$. Therefore $c_2 = \binom{2}{1}_q^2 = (q+1)^2 \equiv 1 \pmod 2$, violating condition (iii) of Theorem~\ref{thm:drg-main}. Thus $J_q(n,k)$ never yields a binary even LCD code when $q$ is even.
\end{proof}

\end{document}
